\subsection{The nilpotent four-element corner on at most four states}

\begin{proposition}[Proved: a nilpotent four-corner small-state reduction]
\label{CP2:A-nilpotent-four-corner-four-states}
Let $N=\{1,e_0,e_1,a\}$ have products
\[
 e_i e_j=e_j,\qquad ae_i=e_i,\qquad e_i a=e_0,
 \qquad a^2=e_0.
\]
Let $S$ be a nonempty left-$N$, right-$M$ biset with $|S|\leq4$.
If $|M|\leq10$, its collage cannot have the rational homology and
cohomology ring of $\CP^2$.
\end{proposition}
\begin{proof}
Write $P_0,P_1,F$ for the left actions of $e_0,e_1,a$ on $S$.
They commute with the right $M$-action and satisfy
\[
 P_iP_j=P_j,\qquad FP_i=P_i,\qquad P_iF=P_0,
 \qquad F^2=P_0.
\]
The images of $P_0$ and $P_1$ are the same nonempty subset
$T\subseteq S$.  Every element of $N$ fixes $T$ pointwise, and $T$
is $M$-invariant.

The corner $BN$ is rationally acyclic, by the one-row minimal-ideal
calculation.  Collapsing this actual corner therefore induces a
rational homology and ring isomorphism from the collage to
\[
 Y=\operatorname{hocofib}\bigl(B(N,S,M)\longrightarrow BM\bigr).
\]
If $|T|=1$, its unique state is fixed by both monoids.  It gives an
actual section of the bridge projection to $BM$.  Consequently that
projection is injective on rational cohomology, and the cofiber exact
sequence makes every degree-two class of $Y$ restrict to zero on
$BM$.  All such classes lift from the suspension quotient $Y/BM$,
so their products vanish.  Hence suppose $|T|\geq2$.

Here is the full natural chain comparison for the remaining cases.
Put $B=\mathbb QN$, $e=e_0$, $f=1-e$ and $z=a-e$.
The right ideals $eB$ and $fB$ are projective, with bases
\[
 eB=\mathbb Q\{e_0,e_1\},\qquad
 fB=\mathbb Q\{f,z\}.
\]
There is an exact right-projective resolution
\[
 \cdots\xrightarrow{\ z\ }fB
 \xrightarrow{\ z\ }fB
 \xrightarrow{\ e_1-e_0\ }eB
 \xrightarrow{\ \varepsilon\ }\mathbb Q\longrightarrow0.
\]
The displayed labels specify the image of the generating idempotent:
the first maps send $f$ to $z$, and the final nonaugmentation map
sends $f$ to $e_1-e_0$.  Indeed $z^2=0$, $zf=z$,
$(e_1-e_0)a=0$, and each successive kernel and image is the indicated
one-dimensional span.  These identities prove exactness in every
degree, rather than only a finite prefix.

Tensoring with $\mathbb QS$ gives a complex of right
$R=\mathbb QM$-modules, natural for the bridge projection.
In its reduced augmented form it is
\[
 \begin{gathered}
 C_0=L_T:=\ker(\mathbb QT\to\mathbb Q),\qquad
 C_j=D:=(1-P_0)\mathbb QS\quad(j\geq1),\\
 d_1=P_1-P_0,\qquad d_j=H:=F-P_0\quad(j\geq2).
 \end{gathered}
\]
In particular $H^2=0$ and $d_1H=0$.  The actual bridge cofiber has
reduced chains computed by
\[
 (C\otimes_R^{\mathbb L}\mathbb Q)[1].
\]
The comparison respects the inclusion of the actual sub-biset $T$:
its reduced augmented complex is precisely $L_T$ in degree zero,
and its inclusion is the displayed inclusion $L_T\to C$.

Since $\dim D=|S|-|T|\leq2$, either $H=0$, or $\dim D=2$ and
$H$ has rank one.  First suppose $H=0$.  As a complex of right
$R$-modules there is an actual direct-sum decomposition
\[
 C=[D\xrightarrow{d_1}L_T]\ \oplus\
       \bigoplus_{j\geq2}D[j].
\]
Derived tensoring preserves this bounded-below direct sum.  If any
$\operatorname{Tor}^R_i(D,\mathbb Q)$ were nonzero, the displayed
copies would give nonzero cofiber homology in arbitrarily high
degrees.  A putative rational projective plane therefore forces
$D\otimes_R^{\mathbb L}\mathbb Q$ to be acyclic.  The inclusion
$L_T\to C$ then becomes a quasi-isomorphism after derived tensoring.
Thus the actual subcollage on $T$ includes by a rational homology
equivalence, and hence induces a rational cohomology ring isomorphism.

On this subcollage $N$ acts trivially on $T$.  Replacing $N$ by the
trivial corner identifies its bridge projection with
$BN\times B(T,M)\to BM$.  As $BN$ is rationally acyclic, this
actual quotient induces a rational homology and ring isomorphism to
the ordinary $M$-act cone on $T$.  It has at most four states, so
Proposition~\ref{CP2:A-ten-six-state-rational} excludes its having
the rational projective-plane homology and ring.

Finally suppose $H$ has rank one.  Then $|T|=2$, $|S|=4$, and
$\ker H=\operatorname{im}H$.  The natural quotient of $C$ is a
quasi-isomorphism to the two-term $R$-complex
\[
 [\,D/\operatorname{im}H\xrightarrow{\overline d_1}L_T\,],
\]
whose two terms are one-dimensional.  If $\overline d_1\ne0$,
it is an isomorphism of $R$-modules, so the derived cofiber is
rationally acyclic.  If $\overline d_1=0$, then $d_1=0$, hence
$P_1=P_0$ on all of $\mathbb QS$ and therefore as functions on $S$.
The map $P_0:S\to T$ now commutes with $P_1,F$, and every right
actor, so it is a proper biset retraction.

The biset-retraction argument also preserves a rational
projective-plane homology ring.  Namely, its actual image subcollage
is a retract, and the relative quotient is a suspension.  If its
second cohomology were zero, the generator of the original second
cohomology would lift from that suspension and have square zero.
Otherwise the retraction and naturality force its second and fourth
cohomology to be one-dimensional with nonzero square, and all other
positive rational homology to vanish.  Thus the image subcollage
would again have the rational projective-plane homology and ring.
Its left $N$-action is trivial on $T$, so the preceding ordinary
act-cone comparison, now on two states, excludes it.  This exhausts
the cases.
\end{proof}

\begin{remark*}[Precise use in the remaining scope]
This proposition removes the unique nilpotent opposite corner from
the $(10,4,4)$ budget without any condition on its source algebra.
It does not by itself remove the five other nonprojective
four-element opposite corners.  For a source carrying a verified
projective or local row-character witness, those five are separately
covered by Theorem~\ref{CP2:A-projective-local-row-graph-obstruction}.
Any remaining source--opposite pieces still require their own proved
reduction or complete finite scope.
\end{remark*}
